\documentclass[11pt]{article}
\usepackage[T1]{fontenc}
\usepackage{lmodern,amsmath,amssymb,amsthm,geometry,microtype,xcolor}
\geometry{margin=1in}
\usepackage[colorlinks=true,allcolors=blue!45!black]{hyperref}
\newtheorem{theorem}{Theorem}
\newtheorem{remark}[theorem]{Remark}
\title{A normalization obstruction to the premagic\newline eigenvector norm conjecture}
\author{Myeongjin Kang}
\date{October 25, 2026}
\begin{document}
\maketitle
\begin{abstract}
Conjecture 7.1 of Teknomo's \emph{Premagic and Ideal Flow Matrices}
asserts that a dominant eigenvector of a nonnegative premagic matrix
has both Euclidean and Frobenius norm one. As stated, the claim is false:
the strictly positive two-by-two all-ones matrix has dominant eigenvector
$(2,2)^T$, whose squared norm is eight. We give a complete Lean certificate,
including a bound on every complex eigenvalue. The issue is the absence of
a normalization convention, rather than a special spectral phenomenon.
\end{abstract}
\noindent\textit{Scope.} This is an elementary correction to a precisely
stated conjecture, not a claim of a new general spectral technique or
worldwide priority. The paper is AI-assisted; original prose is CC BY 4.0
and original code is MIT.

\section{Statement and refutation}
A real square matrix is \emph{premagic} when, for every index, its row sum
equals the corresponding column sum. Teknomo's Definition 4.3 defines this
condition, and Conjecture 7.1 asserts unit Euclidean and Frobenius norms
for the eigenvector associated with the dominant eigenvalue of a
nonnegative premagic matrix~\cite{Teknomo}. The same claim is recorded as
MathDB problem 336487~\cite{MathDB}. No normalization is imposed in the
conjecture's stated hypotheses.

\begin{theorem}
Let
\[
 J=\begin{pmatrix}1&1\\1&1\end{pmatrix},\qquad
 v=\begin{pmatrix}2\\2\end{pmatrix}.
\]
The matrix $J$ is strictly positive and premagic. The vector $v$ is a
nonzero eigenvector with eigenvalue $2$, which has maximal modulus among
all complex eigenvalues of $J$. Nevertheless,
\[
 \|v\|_2^2=\|v\|_F^2=8,
 \qquad\text{so}\qquad \|v\|_2\ne1,\quad \|v\|_F\ne1.
\]
Thus Conjecture 7.1 is false as stated.
\end{theorem}
\begin{proof}
All entries of $J$ are positive, and all row and column sums are $2$.
Direct multiplication gives $Jv=(4,4)^T=2v$. To check dominance over the
entire complex spectrum, suppose $J(a,b)^T=\lambda(a,b)^T$ with
$(a,b)\ne(0,0)$. The coordinate equations are
\[
 a+b=\lambda a,\qquad a+b=\lambda b.
\]
If $a+b=0$, these equations and the nonzero-vector condition imply
$\lambda=0$. If $a+b\ne0$, adding the equations gives
$2(a+b)=\lambda(a+b)$, whence $\lambda=2$. Every eigenvalue therefore has
modulus at most $2$, and the displayed eigenpair is dominant. Finally,
the Euclidean norm and the Frobenius norm of a real column vector are
both the square root of the sum of its squared entries. Here that sum is
$2^2+2^2=8$, which is different from $1$.
\end{proof}

\section{The normalization issue}
If $Av=\lambda v$ and $v\ne0$, then $A(cv)=\lambda(cv)$ for every nonzero
scalar $c$. Consequently an unnormalized eigenvector cannot have its norm
fixed by a property of the matrix alone. Choosing a unit eigenvector is
always possible in a normed vector space, but that is a choice of scale,
not a property distinguishing premagic matrices. If the source intended
an eigenvector already normalized to Euclidean norm one, then its norm
assertion is tautological. The theorem above addresses the literal
unnormalized statement and does not claim that a normalized version is
false.

\section{Formal statement and verification}
The Lean development uses actual mathlib Euclidean space for the norm.
It separately defines the Frobenius norm by
$\sqrt{\sum_i v_i^2}$, proves that the squared Euclidean norm is $8$,
and proves that neither norm is one. Eigenvectors are explicitly required
to be nonzero. The spectral calculation is proved over every field, and
then specialized to the complex numbers to certify the dominance
condition. Positivity and the premagic equations are also checked.

The main declaration is
\texttt{NormOneCounterexample.main\_counterexample}; it combines all of
these conditions. It does not take the eigenvalue bound as an unproved
assumption. Lean 4.34.1 and the pinned mathlib revision from the campaign
are used; the audit records every local theorem's type and transitive
axioms. There are no proof placeholders or added axioms. Source snapshots, hashes, build instructions, and the exact extent of the formal result
accompany the certificate.

This outcome is independent of the separate diagonalizability question
in Conjecture 7.2: the matrix $J$ here is diagonalizable, while its
unnormalized dominant eigenvector still refutes Conjecture 7.1. The present note concerns the distinct assertion about eigenvector norms.

\begin{thebibliography}{9}
\bibitem{Teknomo} K.~Teknomo, \emph{Premagic and Ideal Flow Matrices},
arXiv:1706.08856v2, May 6, 2018, Definition 4.3 and Conjecture 7.1.
\url{https://arxiv.org/abs/1706.08856v2}.
\bibitem{MathDB} MathDB, \emph{Norm-one conjecture for dominant eigenvectors
of non-negative premagic matrices}, problem 336487. Accessed October 25,
2026. \url{https://mathdb.com/p/336487}.
\end{thebibliography}
\end{document}
